\chapter{The Marsh}

\lettrine{A}{thos} and Monk passed over, in going from the camp  towards the Tweed, that part of the ground which Digby had traversed with the  fishermen coming from the Tweed to the camp. The aspect of this place, the  aspect of the changes man had wrought in it, was of a nature to produce a great  effect upon a lively and delicate imagination like that of Athos. Athos looked  at nothing but these desolate spots; Monk looked at nothing but Athos—at  Athos, who, with his eyes sometimes directed towards heaven, and sometimes  towards the earth, sought, thought, and sighed.

Digby, whom the last orders of the general, and particularly the accent with  which he had given them, had at first a little excited, Digby followed the pair  at about twenty paces, but the general having turned round as if astonished to  find his orders had not been obeyed, the aid-de-camp perceived his  indiscretion, and returned to his tent.

He supposed that the general wished to make, incognito, one of those reviews of  vigilance which every experienced captain never fails to make on the eve of a  decisive engagement: he explained to himself the presence of Athos in this case  as an inferior explains all that is mysterious on the part of his leader. Athos  might be, and, indeed, in the eyes of Digby, must be, a spy, whose information  was to enlighten the general.

At the end of a walk of about ten minutes among the tents and posts, which were  closer together near the headquarters, Monk entered upon a little causeway  which diverged into three branches. That on the left led to the river, that in  the middle to Newcastle Abbey on the marsh, that on the right crossed the first  lines of Monk's camp; that is to say, the lines nearest to  Lambert's army. Beyond the river was an advanced post, belonging to  Monk's army, which watched the enemy; it was composed of one hundred and  fifty Scots. They had swum across the Tweed, and, in case of attack, were to  recross it in the same manner, giving the alarm; but as there was no post at  that spot, and as Lambert's soldiers were not so prompt at taking to the  water as Monk's were, the latter appeared not to have as much uneasiness  on that side. On this side of the river, at about five hundred paces from the  old Abbey, the fishermen had taken up their abode amidst a crowd of small tents  raised by soldiers of the neighbouring clans, who had with them their wives and  children. All this confusion, seen by the moon's light, presented a  striking coup d'œil; the half shadow enlarged every detail, and the  light, that flatterer which only attaches itself to the polished side of  things, courted upon each rusty musket the point still left intact, and upon  every rag of canvas the whitest and least sullied part.

Monk arrived then with Athos, crossing this spot, illumined with a double  light, the silver splendour of the moon, and the red blaze of the fires at the  meeting of these three causeways; there he stopped, and addressing his  companion,—<Monsieur,> said he, <do you know your  road?>

<General, if I am not mistaken, the middle causeway leads straight to the  Abbey.>

<That is right; but we shall want lights to guide us in the  vaults.> Monk turned round.

<Ah! I thought Digby was following us!> said he. <So much the  better; he will procure us what we want.>

<Yes, general, there is a man yonder who has been walking behind us for  some time.>

<Digby!> cried Monk. <Digby! come here, if you please.>

But instead of obeying, the shadow made a motion of surprise, and, retreating  instead of advancing, it bent down and disappeared along the jetty on the left,  directing its course towards the lodging of the fishermen.

<It appears not to be Digby,> said Monk.

Both had followed the shadow which had vanished. But it was not so rare a thing  for a man to be wandering about at eleven o'clock at night, in a camp in  which are reposing ten or eleven thousand men, as to give Monk and Athos any  alarm at his disappearance.

<As it is so,> said Monk, <and we must have a light, a  lantern, a torch, or something by which we may see where to see our feet; let  us seek this light.>

<General, the first soldier we meet will light us.>

<No,> said Monk, in order to discover if there were not any  connivance between the Comte de la Fère and the fisherman. <No, I should  prefer one of these French sailors who came this evening to sell me their fish.  They leave to-morrow, and the secret will be better kept by them; whereas, if a  report should be spread in the Scottish army, that treasures are to be found in  the Abbey of Newcastle, my Highlanders will believe there is a million  concealed beneath every slab, and they will not leave stone upon stone in the  building.>

<Do as you think best, general,> replied Athos, in a natural tone  of voice, making evident that soldier or fisherman was the same to him, and  that he had no preference.

Monk approached the causeway behind which had disappeared the person he had  taken for Digby, and met a patrol who, making the tour of the tents, was going  towards headquarters; he was stopped with his companion, gave the password, and  went on. A soldier, roused by the noise, unrolled his plaid, and looked up to  see what was going forward. <Ask him,> said Monk to Athos,  <where the fishermen are; if I were to speak to him, he would know  me.>

Athos went up to the soldier, who pointed out the tent to him; immediately Monk  and Athos turned towards it. It appeared to the general that at the moment they  came up, a shadow like that they had already seen, glided into this tent; but  on drawing nearer he perceived he must have been mistaken, for all of them were  asleep pele mele, and nothing was seen but arms and legs joined, crossed, and  mixed. Athos, fearing lest he should be suspected of connivance with some of  his compatriots, remained outside the tent.

<Hola!> said Monk, in French, <wake up here.> Two or  three of the sleepers got up.

<I want a man to light me,> continued Monk.

<Your honour may depend on us,> said a voice which made Athos start.  <Where do you wish us to go?>

<You shall see. A light! come, quickly!>

<Yes, your honour. Does it please your honour that I should accompany  you?>

<You or another; it is of very little consequence, provided I have a  light.>

<It is strange!> thought Athos; <what a singular voice that  man has!>

<Some fire, you fellows!> cried the fisherman; <come, make  haste!>

Then addressing his companion nearest to him in a low voice:—<Get  ready a light, Menneville,> said he, <and hold yourself ready for  anything.>

One of the fishermen struck light from a stone, set fire to some tinder, and by  the aid of a match lit a lantern. The light immediately spread all over the  tent.

<Are you ready, monsieur?> said Monk to Athos, who had turned away,  not to expose his face to the light.

<Yes, general,> replied he.

<Ah! the French gentleman!> said the leader of the fishermen to  himself. <Peste! I have a great mind to charge you with the commission,  Menneville; he may know me. Light! light!> This dialogue was pronounced  at the back of the tent, and in so low a voice that Monk could not hear a  syllable of it; he was, besides, talking with Athos. Menneville got himself  ready in the meantime, or rather received the orders of his leader.

<Well?> said Monk.

<I am ready, general,> said the fisherman.

Monk, Athos, and the fisherman left the tent.

<It is impossible!> thought Athos. <What dream could put that  into my head?>

<Go forward; follow the middle causeway, and stretch out your  legs,> said Monk to the fisherman.

They were not twenty paces on their way when the same shadow that had appeared  to enter the tent came out of it again, crawled along as far as the piles, and,  protected by that sort of parapet placed along the causeway, carefully observed  the march of the general. All three disappeared in the night haze. They were  walking towards Newcastle, the white stones of which appeared to them like  sepulchres. After standing for a few seconds under the porch, they penetrated  into the interior. The door had been broken open by hatchets. A post of four  men slept in safety in a corner, so certain were they that the attack would not  take place on that side.

<Will not these men be in your way?> said Monk to Athos.

<On the contrary, monsieur, they will assist in rolling out the barrels,  if your honour will permit them.>

<You are right.>

The post, though fast asleep, roused up at the first steps of the three  visitors amongst the briars and grass that invaded the porch. Monk gave the  password, and penetrated into the interior of the convent, preceded by the  light. He walked last, watching the least movement of Athos, his naked dirk in  his sleeve, and ready to plunge it into the back of the gentleman at the first  suspicious gesture he should see him make. But Athos, with a firm and sure  step, crossed the chambers and courts.

Not a door, not a window was left in the building. The doors had been burnt,  some on the spot, and the charcoal of them was still jagged with the action of  the fire, which had gone out of itself, powerless, no doubt, to get to the  heart of those massive joints of oak fastened together with iron nails. As to  the windows, all the panes having been broken, night birds, alarmed by the  torch, flew away through their holes. At the same time, gigantic bats began to  trace their vast, silent circles around the intruders, whilst the light of the  torch made their shadows tremble on the high stone walls. Monk concluded that  there could be no man in the convent, since wild beasts and birds were there  still, and fled away at his approach.

After having passed the rubbish, and torn away more than one branch of ivy that  had made itself a guardian of the solitude, Athos arrived at the vaults  situated beneath the great hall, but the entrance of which was from the chapel.  There he stopped. <Here we are, general,> said he.

<This, then, is the slab?>

\begin{bwbigpic}
	[\basicscale] 
	{25slab}
	[<This, then, is the slab?>]
\end{bwbigpic}

<Yes.>

<Ay, and here is the ring—but the ring is sealed into the  stone.>

<We must have a lever.>

<That's a very easy thing to find.>

Whilst looking around them, Athos and Monk perceived a little ash of about  three inches in diameter, which had shot up in an angle of the wall, reaching a  window, concealed by its branches.

<Have you a knife?> said Monk to the fisherman.

<Yes, monsieur.>

<Cut down this tree, then.>

The fisherman obeyed, but not without notching his cutlass. When the ash was  cut and fashioned into the shape of a lever, the three men penetrated into the  vault.

<Stop where you are,> said Monk to the fisherman. <We are  going to dig up some powder; your light may be dangerous.>

The man drew back in a sort of terror, and faithfully kept to the post assigned  him, whilst Monk and Athos turned behind a column at the foot of which,  penetrating through a crack, was a moonbeam, reflected exactly on the stone  which the Comte de la Fère had come so far in search.

<This is it,> said Athos, pointing out to the general the Latin  inscription.

<Yes,> said Monk.

Then, as if still willing to leave the Frenchman one means of evasion,—

<Do you not observe that this vault has already been broken into,>  continued he, <and that several statues have already been knocked  down?>

<My lord, you have, without doubt, heard that the religious respect of  your Scots loves to confide to the statues of the dead the valuable objects  they have possessed during their lives. Therefore, the soldiers had reason to  think that under the pedestals of the statues which ornament most of these  tombs, a treasure was hidden. They have consequently broken down pedestal and  statue: but the tomb of the venerable cannon, with which we have to do, is not  distinguished by any monument. It is simple, therefore it has been protected by  the superstitious fear which your Puritans have always had of sacrilege. Not a  morsel of the masonry of this tomb has been chipped off.>

<That is true,> said Monk.

Athos seized the lever.

<Shall I help you?> said Monk.

<Thank you, my lord; but I am not willing that your honour should lend  your hand to a work of which, perhaps, you would not take the responsibility if  you knew the probable consequences of it.>

Monk raised his head.

<What do you mean by that, monsieur?>

<I mean—but that man\longdash>

<Stop,> said Monk; <I perceive what you are afraid of. I  shall make a trial.> Monk turned towards the fisherman, the whole of  whose profile was thrown upon the wall.

<Come here, friend!> said he in English, and in a tone of command.

The fisherman did not stir.

<That is well,> continued he: <he does not know English.  Speak to me, then, in English, if you please, monsieur.>

<My lord,> replied Athos, <I have frequently seen men in  certain circumstances have sufficient command over themselves not to reply to a  question put to them in a language they understood. The fisherman is perhaps  more learned than we believe him to be. Send him away, my lord, I beg  you.>

<Decidedly,> said Monk, <he wishes to have me alone in this  vault. Never mind, we shall go through with it; one man is as good as another  man; and we are alone. My friend,> said Monk to the fisherman, <go  back up the stairs we have just descended, and watch that nobody comes to  disturb us.> The fisherman made a sign of obedience. <Leave your  torch,> said Monk; <it would betray your presence, and might  procure you a musket-ball.>

The fisherman appeared to appreciate the counsel; he laid down the light, and  disappeared under the vault of the stairs. Monk took up the torch, and brought  it to the foot of the column.

<Ah, ah!> said he; <money, then, is concealed under this  tomb?>

<Yes, my lord; and in five minutes you will no longer doubt it.>

At the same time Athos struck a violent blow upon the plaster, which split,  presenting a chink for the point of the lever. Athos introduced the bar into  this crack, and soon large pieces of plaster yielded, rising up like rounded  slabs. Then the Comte de la Fère seized the stones and threw them away with a  force that hands so delicate as his might not have been supposed capable of  having.

<My lord,> said Athos, <this is plainly the masonry of which  I told your honour.>

<Yes; but I do not yet see the casks,> said Monk.

<If I had a dagger,> said Athos, looking round him, <you  should soon see them, monsieur. Unfortunately, I left mine in your tent.>

<I would willingly offer you mine,> said Monk, <but the blade  is too thin for such work.>

Athos appeared to look around him for a thing of some kind that might serve as  a substitute for the weapon he desired. Monk did not lose one of the movements  of his hands, or one of the expressions of his eyes. <Why do you not ask  the fisherman for his cutlass?> said Monk; <he has a  cutlass.>

<Ah! that is true,> said Athos; <for he cut the tree down  with it.> And he advanced towards the stairs.

<Friend,> said he to the fisherman, <throw me down your  cutlass, if you please; I want it.>

The noise of the falling weapon sounded on the steps.

<Take it,> said Monk; <it is a solid instrument, as I have  seen, and a strong hand might make good use of it.>

Athos appeared only to give to the words of Monk the natural and simple sense  under which they were to be heard and understood. Nor did he remark, or at  least appear to remark, that when he returned with the weapon, Monk drew back,  placing his left hand on the stock of his pistol; in the right he already held  his dirk. He went to work then, turning his back to Monk, placing his life in  his hands, without possible defence. He then struck, during several seconds, so  skilfully and sharply upon the intermediary plaster, that it separated into  two parts, and Monk was able to discern two barrels placed end to end, and  which their weight maintained motionless in their chalky envelope.

<My lord,> said Athos, <you see that my presentiments have  not been disappointed.>

<Yes, monsieur,> said Monk, <and I have good reason to  believe you are satisfied; are you not?>

<Doubtless I am; the loss of this money would have been inexpressibly  great to me: but I was certain that God, who protects the good cause, would not  have permitted this gold, which should procure its triumph, to be diverted to  baser purposes.>

<You are, upon my honour, as mysterious in your words as in your actions,  monsieur,> said Monk. <Just now as I did not perfectly understand  you when you said that you were not willing to throw upon me the responsibility  of the work we were accomplishing.>

<I had reason to say so, my lord.>

<And now you speak to me of the good cause. What do you mean by the words  <the good cause?> We are defending at this moment, in England, five  or six causes, which does not prevent every one from considering his own not  only as the good cause, but as the best. What is yours, monsieur? Speak boldly,  that we may see if, upon this point, to which you appear to attach a great  importance, we are of the same opinion.>

Athos fixed upon Monk one of those penetrating looks which seemed to convey to  him to whom they are directed a challenge to conceal a single one of his  thoughts; then, taking off his hat, he began in a solemn voice, while his  interlocutor, with one hand upon his visage, allowed that long and nervous hand  to compress his moustache and beard, while his vague and melancholy eye wandered  about the recesses of the vaults.  